% -*- coding: utf-8 -*-
% !TEX program = xelatex
\documentclass[plain,noanswer,medmath]{../../template/jnuexam}
\usepackage{../../template/kysx}

\begin{document}


\examtitle{2005}{3}


\loadproblems{1}{2005P1}%载入试卷一
\loadproblems{2}{2005P2}%载入试卷二


\makepart{填空题}{1～6小题，每小题4分，共24分}


\begin{problem}
极限$\lim_{x \to \infty} x\sin \frac{2x}{x^2 + 1} =$ \fillin{}．
\end{problem}


\begin{problem}
微分方程$xy' + y = 0$满足初始条件$y(1) = 2$的特解为 \fillin{}．
\end{problem}


\begin{problem}
设二元函数$z = x\e^{x + y} + (x + 1)\ln(1 + y)$，则$\dz\big|_{(1,0)} =$ \fillin{}．
\end{problem}


\begin{problem}
设行向量组$(2,1,1,1)$，$(2,1,a,a)$，$(3,2,1,a)$，$(4,3,2,1)$线性相关，且$a \ne 1$，则$a =$ \fillin{}．
\end{problem}


\useproblem{1}{1}{6}%【同试卷一第一(6)题】


%【类似试卷一第二(13)题】
\begin{problem}
设二维随机变量$(X,Y)$的概率分布为
$$\begin{array}{|c|cc|}
\hline
  \diagboxtwo{X}{Y} &   0 & 1 \\
\hline
                  0 & 0.4 & a \\
                  1 &   b & 0.1 \\
\hline
\end{array}$$
已知随机事件$\{X = 0\}$与$\{X + Y = 1\}$相互独立，则$a=$ \fillin{}，$b=$ \fillin{}．
\end{problem}


\makepart{选择题}{7～14小题，每小题4分，共32分}


\begin{problem}
当$a$取下列哪个值时，函数$f(x) = 2x^3 - 9x^2 + 12x - a$恰好有两个不同的零点\pickout{}
\begin{abcd}
\item $2$．
\item $4$．
\item $6$．
\item $8$．
\end{abcd}
\end{problem}


\begin{problem}
设$I_1 = \iint_D \cos \sqrt{x^2 + y^2} \d\sigma$, $I_2 = \iint_D \cos(x^2 + y^2)\d\sigma$,
$I_3 = \iint_D \cos(x^2 + y^2)^2 \d\sigma$，其中
$D = \left\{(x,y) \middle| x^2 + y^2 \le 1 \right\}$，则\pickout{}
\begin{abcd}
\item $I_3 > I_2 > I_1$．
\item $I_1 > I_2 > I_3$．
\item $I_2 > I_1 > I_3$．
\item $I_3 > I_1 > I_2$．
\end{abcd}
\end{problem}


\begin{problem}
设$a_n > 0$, $n = 1,2, \cdots$，若$\sum_{n = 1}^{\infty}a_n$发散，
$\sum_{n = 1}^{\infty}(- 1)^{n - 1}a_n$收敛，则下列结论正确的是\pickout{}
\begin{abcd}
\item $\sum_{n = 1}^{\infty}a_{2n - 1}$收敛，$\sum_{n = 1}^{\infty}a_{2n}$发散．
\item $\sum_{n = 1}^{\infty}a_{2n}$收敛，$\sum_{n = 1}^{\infty}a_{2n - 1}$发散．
\item $\sum_{n = 1}^{\infty}(a_{2n - 1} + a_{2n})$收敛．
\item $\sum_{n = 1}^{\infty}(a_{2n - 1} - a_{2n})$收敛．
\end{abcd}
\end{problem}


\begin{problem}
设$f(x) = x\sin x + \cos x$，下列命题中正确的是\pickout{}
\begin{abcd}
\item $f(0)$是极大值，$f(\frac{\pi}{2})$是极小值．
\item $f(0)$是极小值，$f(\frac{\pi}{2})$是极大值．
\item $f(0)$是极大值，$f(\frac{\pi}{2})$也是极大值．
\item $f(0)$是极小值，$f(\frac{\pi}{2})$也是极小值．
\end{abcd}
\end{problem}


\begin{problem}
以下四个命题中，正确的是\pickout{}
\begin{abcd}
\item 若$f'(x)$在$(0,1)$内连续，则$f(x)$在$(0,1)$内有界．
\item 若$f(x)$在$(0,1)$内连续，则$f(x)$在$(0,1)$内有界．
\item 若$f'(x)$在$(0,1)$内有界，则$f(x)$在$(0,1)$内有界．
\item 若$f(x)$在$(0,1)$内有界，则$f'(x)$在$(0,1)$内有界．
\end{abcd}
\end{problem}


\begin{problem}
设矩阵$A$=$(a_{ij})_{3 \times 3}$满足$A^* = A^T$，其中$A^*$是$A$的伴随矩阵，$A^T$为$A$的转置矩阵．
若$a_{11},a_{12},a_{13}$为三个相等的正数，则$a_{11}$为\pickout{}
\begin{abcd}
\item $\frac{\sqrt{3}}{3}$．
\item $3$．
\item $\frac{1}{3}$．
\item $\sqrt{3}$．
\end{abcd}
\end{problem}


\useproblem{1}{2}{11}%【同试卷一第二(11)题】


\begin{problem}
设一批零件的长度服从正态分布$N(\mu ,\sigma^2)$，其中$\mu ,\sigma^2$均未知．
现从中随机抽取$16$个零件，测得样本均值$\overline{x} = 20(cm)$，样本标准差$s = 1(cm)$，
则$\mu$的置信度为$0.90$的置信区间是\pickout{}
\begin{abcd}
\item $\left(20 - \frac{1}{4}t_{0.05}(16),20 + \frac{1}{4}t_{0.05}(16) \right)$．
\item $\left(20 - \frac{1}{4}t_{0.1}(16),20 + \frac{1}{4}t_{0.1}(16) \right)$．
\item $\left(20 - \frac{1}{4}t_{0.05}(15),20 + \frac{1}{4}t_{0.05}(15) \right)$．
\item $\left(20 - \frac{1}{4}t_{0.1}(15),20 + \frac{1}{4}t_{0.1}(15) \right)$．
\end{abcd}
\end{problem}


\makepart{解答题}{本题共9小题，满分94分}


\begin{problem}[points=8]
求$\lim_{x \to 0} \left(\frac{1 + x}{1 - \e^{- x}} - \frac{1}{x} \right)$．
\end{problem}


\begin{problem}[points=8]
设$f(u)$具有二阶连续导数，且$g(x,y) = f\left(\frac{y}{x} \right) + yf\left(\frac{x}{y} \right)$，
求$x^2\frac{\pd^2 g}{\pd x^2} - y^2\frac{\pd^2 g}{\pd y^2}$．
\end{problem}


\useproblem[points=9]{2}{3}{21}%【同试卷二第三(21)题】


\begin{problem}[points=9]
求幂级数$\sum_{n = 1}^{\infty}\left(\frac{1}{2n + 1} - 1 \right) x^{2n}$在区间$(- 1,1)$内的和函数$S(x)$．
\end{problem}


\begin{problem}[points=8]
设$f(x),g(x)$在$[0,1]$上的导数连续，且$f(0) = 0$,$f'(x) \ge 0$,$g'(x) \ge 0$．
证明：对任何$a \in[0,1]$，有$\int_0^a g(x)f'(x)\dx + \int_0^1 f(x)g'(x)\dx \ge f(a)g(1).$
\end{problem}


\begin{problem}[points=13]
已知齐次线性方程组
(I)$\begin{cases}
  x_1 + 2x_2 + 3x_3 = 0, \\
  2x_1 + 3x_2 + 5x_3 = 0, \\
  x_1 + x_2 + ax_3 = 0,
\end{cases}$和(II)$\begin{cases}
  x_1 + bx_2 + cx_3 = 0, \\
  2x_1 + b^2x_2 + (c + 1)x_3 = 0,
\end{cases}$
同解，求$a,b,c$的值．
\end{problem}


\begin{problem}[points=13]
设$D = \begin{pmatrix}
    A & C \\
  C^T & B
\end{pmatrix}$为正定矩阵，其中$A,B$分别为$m$阶，$n$阶对称矩阵，$C$为$m \times n$矩阵．
\begin{enumerate}
\item 计算$P^T DP$，其中$P = \begin{pmatrix}
  E_m & -A^{-1} C \\
    O & E_n
\end{pmatrix}$；
\item 利用(I)的结果判断矩阵$B - C^T A^{-1} C$是否为正定矩阵，并证明你的结论．
\end{enumerate}
\end{problem}


%【类似试卷一第三(22)题】
\begin{problem}[points=13]
设二维随机变量$(X,Y)$的概率密度为$f(x,y) = \begin{cases}
  1, & 0 < x < 1,0 < y < 2x, \\
  0, & \text{其他}.
\end{cases}$求：
\begin{enumerate}
\item $(X,Y)$的边缘概率密度$f_X(x),f_Y(y)$；
\item $Z = 2X - Y$的概率密度$f_Z(z)$；
\item $P\left\{Y \le \frac{1}{2} \middle| X \le \frac{1}{2} \right\}$．
\end{enumerate}
\end{problem}


%【类似试卷一第三(23)题】
\begin{problem}[points=13]
设$X_1,X_2, \cdots ,X_n(n > 2)$为来自总体$N(0,\sigma^2)$的简单随机样本，
其样本均值为$\overline{X}$，记$Y_i = X_i - \overline{X}$，$i = 1,2, \cdots ,n$．
\begin{enumerate}
\item 求$Y_i$的方差$DY_i$，$i = 1,2, \cdots ,n$；
\item 求$Y_1$与$Y_n$的协方差$\Cov (Y_1,Y_n)$；
\item 若$c(Y_1 + Y_n)^2$是$\sigma^2$的无偏估计量，求常数$c$．
\end{enumerate}
\end{problem}


\end{document}
